\documentclass[10pt,reqno]{amsart}
\usepackage[T1]{fontenc}
\usepackage{lmodern,microtype}
\usepackage{amsmath,amssymb,mathtools}
\usepackage[colorlinks=true,linkcolor=blue,citecolor=blue,urlcolor=blue]{hyperref}
\newtheorem{theorem}{Theorem}[section]
\newtheorem{lemma}[theorem]{Lemma}
\newtheorem{corollary}[theorem]{Corollary}
\theoremstyle{remark}
\newtheorem{remark}[theorem]{Remark}
\newcommand{\K}{\mathcal K}
\newcommand{\J}{\mathcal J}
\newcommand{\A}{\mathcal A}
\newcommand{\Sset}{\mathcal S}
\newcommand{\E}{\mathbb E}
\newcommand{\Prob}{\mathbb P}
\newcommand{\ind}{\mathbf1}
\newcommand{\li}{\operatorname{li}}
\newcommand{\lcm}{\operatorname{lcm}}
\newcommand{\rad}{\operatorname{rad}}
\newcommand{\e}{\mathrm e}
\title[An internally proved three-quarter exceptional-set bound]{A three-quarter exceptional-set bound for the\ Erd\H{o}s--Straus equation: an internally proved research note}
\author{Anonymous}
\date{INTERNALLY PROVED; internal checks only, not externally refereed}
\begin{document}
\begin{abstract}
This note, \textbf{internally proved but not externally refereed},
records the internally reviewed argument for
\[
 E(N)\ll N\exp\{-c(\log N)^{3/4}\},
\]
where $E(N)$ counts denominators
$n\leq N$ for which $4/n$ is not a sum of three positive unit fractions.
Power-sized multipliers supply cubic logarithmic class mass after an
incidence-congestion pruning and Bombieri--Vinogradov.  A fixed family of
congruence classes is then assembled on the integers, using exact CRT
probabilities, a small-prime selector, and an even Bonferroni polynomial.
The complete coefficient and rounding ledger costs $\exp\{O((\log X)^4)\}$.
All proof ingredients and their uniformities are given below; the three
external-referee checkpoints are retained explicitly, and
Section~\ref{sec:unpruned} gives a shorter route to the class supply
(Cauchy--Schwarz against Bombieri--Vinogradov) that removes the
distinct-prime-factor cutoff, the congestion pruning, and the unweighted
incidence-moment lemma (Lemma~\ref{lem:congestion}) from the proof.  The
argument has undergone the internal checks described in
Theorem~\ref{thm:main}; these do not constitute external refereeing,
and ``internally proved'' means exactly that.  This is not a proof of the
Erd\H{o}s--Straus conjecture, nor a claim of publication priority.
\end{abstract}
\maketitle

\section{Statement, status, and parameter order}\label{sec:intro}
For a positive integer $n$, call $n$ \emph{representable} if there exist
positive integers $x_1,x_2,x_3$ with
\begin{equation}\label{eq:ES}
 \frac4n=\frac1{x_1}+\frac1{x_2}+\frac1{x_3}.
\end{equation}
The denominators $x_i$ need not be distinct.  Put
\[
 \begin{aligned}
 E(N)&=\#\{1\leq n\leq N:n\text{ is not representable}\},\\
 E_{\rm pr}(N)&=\#\{p\leq N:p\text{ prime and not representable}\}.
 \end{aligned}
\]
The contribution of $n=1$ is harmless.  The conjecture asks for no
exceptions for $n\geq2$; a density bound does not establish that assertion.

\begin{theorem}[internally proved; not externally refereed]\label{thm:main}
There are absolute $c_*,C_*>0$ such that, for every real
$N\geq3$,
\[
 E_{\rm pr}(N)\leq E(N)\leq C_*N\exp\{-c_*(\log N)^{3/4}\}.
\]
The argument for this statement has undergone internal checks only.
The internal record comprises a self-review (\cite[\S39.7]{Notes}), a
blind construction and reconciliation taking Theorem~34.8 there as
input (\cite[\S41]{Notes}; the blindness is self-attested, not
independently audited), a standalone internal audit (the repository
file \texttt{reviews/es-threequarter-note-review.md}, which found no
load-bearing gap and recorded repairs), a checkpoint re-derivation
(\cite[\S76.4]{Notes}), and a hostile review of Section~\ref{sec:unpruned}
(the file \texttt{reviews/wave32-sec76-review.md}).  These records are not three
independent full-chain hostile reviews, and the theorem has not been
verified by an external expert.
\end{theorem}

This note follows \cite[\S\S16,34,39]{Notes}; \cite[\S41]{Notes} is an
internal blind re-derivation, not an independent published source.
We use the reviewed weaker-bound note \cite{LL} as follows:
its Lemma~2.1 supplies the multiplier identity; Lemmas~3.1 and~3.2
supply elementary lattice sums and the multiplier harmonic sum;
the distinctness argument is that of its Lemma~4.1; its Section~5 fixes
our large-sieve conventions; and its Section~6 supplies the final
semigroup argument.  These ingredients are restated, and their proofs
are reproduced or lightly adapted, so no consultation of a working
notebook is needed to follow the calculation.  We do \emph{not} import
\cite[Lemma~4.1]{LL}'s prime-counting conclusion outside its stated
polylogarithmic multiplier range.  That extension is the work of
Section~\ref{sec:harvest}.

Here is the order of constants.  Fix once and for all
\begin{equation}\label{eq:parameters}
 \begin{gathered}
 0<\kappa<1/240,\quad t=\log X,\quad
 K=\lfloor X^\kappa\rfloor,\quad H=K^{10},\\
 \K(K)=\{1\leq k\leq K:k\equiv1\pmod4\},\quad
 L_K=\lcm_{k\in\K(K)}k.
 \end{gathered}
\end{equation}
An absolute sufficiently large $K_0$ is chosen in the lattice estimates.
The distinct-prime-factor cutoff $D$ is a fixed absolute integer chosen
in Lemma~\ref{lem:omega}.  Constants in the analytic estimates are
absolute once $D$ is fixed, with thresholds allowed to depend on
$\kappa,D$.  In Section~\ref{sec:void} we choose an absolute small
$\eta>0$, obtain a positive void constant $c_v=c_v(\kappa,D)$, and then
fix a sufficiently large $B=B(\kappa,D)$ for $y=Bt^3$.
The Bonferroni constant $D_B=D_B(\kappa,D,B)$ is distinct from $D$.
Finally choose $C_0=C_0(\kappa,D,B,D_B)$ for the ledger and
$\alpha=\alpha(C_0)>0$ for $t=\alpha(\log N)^{1/4}$.
All subsequent sufficiently-large thresholds have only these dependencies;
no constant depends on a residue, subfamily, moment order, or $N$.
Fixing, for example, $\kappa=1/480$ makes the final constants absolute.
Constants are not asserted to be effective: standard Bombieri--Vinogradov
is used without an effectivity refinement.
This $D$-dependent list describes the retained original route.  For the
simplified route of Section~\ref{sec:unpruned}, omit $D$ and its choice,
use the family $\A'_X$, and delete $D$ from all dependencies; the
remaining choices occur in the same order: $\kappa$, $K_0$, $\eta$ and
the void constant, $B$, $D_B$, $C_0$, and $\alpha$.

The central distinction is between using prime distribution at scale $X$
to \emph{construct} many classes and counting those classes on
\emph{all integers} up to $N$.  A prime-progression error
$N/(\log N)^A$, with fixed $A$, is much larger than the desired
$N\e^{-ct^3}$ in the final window.  Such an error is never used at scale
$N$.  Exact integer counts, each with error at most one, replace it.

\section{Forced classes and a fixed atom family}\label{sec:atoms}
\begin{lemma}[multiplier identity; {\cite[Lemma 2.1]{LL}}]\label{lem:identity}
Let $k,\ell$ be positive integers with $k\ell\equiv3\pmod4$, and let
$A=(k\ell+1)/4=uvw$ be any factorization into positive integers.  If $n\geq1$
satisfies $nv\equiv-u\pmod{k\ell}$, then $s=(nv+u)/(k\ell)$ is a positive
integer and
\[
 \frac4n=\frac1{suw}+\frac1{nsvw}+\frac1{nuvw}.
\]
In particular every positive integer in the class $-uv^{-1}\pmod{k\ell}$ is
representable; the inverse exists since $(A,k\ell)=1$.
\end{lemma}
\begin{proof}
On the common denominator $nsuvw$ the three numerators sum to
$nv+u+s=sk\ell+s=s(k\ell+1)=4suvw$.  Positivity and integrality of $s$
follow directly from the hypotheses.
\end{proof}

For a precise endpoint convention take $x_j=2^jX^{1/2}$ and use the disjoint
blocks
\[
 I_j=(x_j,\min(2x_j,X)],\qquad
 0\leq j<\left\lceil\frac{\log X}{2\log2}\right\rceil.
\]
The last may be partial.  In a block put $z_j=x_j^{1/6}$; we discard the
last block for lower estimates if it is partial and enlarge it to a full
block for upper estimates.  There are $\asymp t$ full blocks and
$\log x_j\asymp t$.  Because $\kappa<1/240$, for all sufficiently large
$X$ each block satisfies
\begin{equation}\label{eq:floor}
 z_j\geq X^{1/12}>K^{20}=H^2,
 \qquad \Lambda_j:=\log(z_j/H)\geq\tfrac12\log z_j.
\end{equation}

Let $\A_X$ consist of all quadruples $A=(k,\ell,u,v)$ such that
\begin{equation}\label{eq:atomdef}
 \begin{gathered}
 k\in\K(K),\quad \ell\in I_j\text{ prime},\quad H<u,v\leq z_j,\\
 (u,v)=(uv,k)=1,\quad \omega(uv)\leq D\log\log X,
 \quad 4uv\mid k\ell+1.
 \end{gathered}
\end{equation}
Here $\omega(m)$ counts distinct prime factors.  Divisibility already forces
$\ell\equiv3\pmod4$.  The atom denotes the cylinder
\begin{equation}\label{eq:cylinder}
 E_A=\{n:n\equiv-uv^{-1}\pmod{k\ell}\},\qquad
 H_X(n)=\sum_{A\in\A_X}\ind_{E_A}(n).
\end{equation}
This is a fixed, unpruned, residue-independent family.  Pruning will be used
only to prove a lower mass bound within a fibre, never to define the
polynomial on the integers.

\begin{lemma}[deduplication and conditional independence]\label{lem:CRT}
For every sufficiently large $X$ in \eqref{eq:parameters} the following hold.
Every exceptional integer avoids every atom.  At a fixed $\ell$ all
atoms have distinct projections modulo $\ell$.  Consequently a compatible
set of distinct atoms has distinct large primes $\ell_i$, and its exact
CRT probability is
\begin{equation}\label{eq:intersection}
 \Prob\left(\bigcap_i E_{A_i}\right)
 =\frac{\ind_{\rm compatible}}
 {\lcm(k_1,\ldots,k_j)\prod_i\ell_i}.
\end{equation}
Conditional on any fixed residue $c\pmod{L_K}$, the coordinates
$n\pmod\ell$ are independent and uniform.  An atom is active in that
fibre precisely when $u+cv\equiv0\pmod k$.
\end{lemma}
\begin{proof}
The identity proves avoidance and reduction modulo $k$ proves the activation
condition.  If two atoms at the same $\ell\in I_j$ have the same projection,
then $\ell\mid uv'-u'v$, but $|uv'-u'v|<z_j^2<\ell$.
Thus $uv'=u'v$, and both fractions being reduced gives $(u,v)=(u',v')$.
Now $k\ell\equiv k'\ell\equiv-1\pmod{4uv}$.  Since $(\ell,4uv)=1$,
we get $k\equiv k'\pmod{4uv}$; as $4uv>4H^2>K$, this means $k=k'$.
In particular there are at most $z_j^2\leq\ell^{1/3}$ atoms at that prime.
This is the distinctness argument of \cite[Lemma~4.1]{LL}.
Finally all $\ell>X^{1/2}>K$ are distinct primes coprime to $L_K$.
The Chinese remainder theorem proves both the density and independence
assertions.  Incompatible intersections are empty.
\end{proof}

All probability spaces below are finite and exact.  When the selector
$S_y$ is present we use the uniform residue $n$ modulo
\begin{equation}\label{eq:space}
 \mathcal M=\lcm(L_K,P_y)\prod_{\substack{X^{1/2}<\ell\leq X\\\ell\ \mathrm{prime}}}\ell,
 \qquad P_y=\prod_{p\leq y}p,\quad 2\leq y<X^{1/2}.
\end{equation}
Without the selector one can take $P_y=1$.  A congruence whose modulus
 divides $\mathcal M$ has probability the reciprocal of its modulus.
The possibly enormous size of $\mathcal M$ is a device for defining
probabilities, not a partition of $[1,N]$ that will be counted term by term.
The factor 24 used in the notebooks may be included in $\mathcal M$ without
changing any argument; no multiplier here requires it.

\section{Analytic inputs and the power-range lattice estimates}\label{sec:lattice}
We record the external inputs precisely.  For $q\geq1$ and $(a,q)=1$,
Brun--Titchmarsh in the Montgomery--Vaughan form \cite{MV} gives
\begin{equation}\label{eq:BT}
 \pi(V;q,a)\leq\frac{2V}{\varphi(q)\log(V/q)}\qquad(V>q).
\end{equation}
Only $q\leq4x^{1/3}$, $V=2x$ is needed.  For every fixed $R>0$,
Bombieri--Vinogradov \cite[Chapter~28]{Davenport} supplies constants
$A_R,C_R,x_R>0$ such that, for all $x\geq x_R$,
\begin{equation}\label{eq:BV}
 \sum_{q\leq x^{1/2}(\log x)^{-A_R}}E_x^*(q)
 \leq C_Rx(\log x)^{-R},
\end{equation}
where
\[
 E_x^*(q)=\max_{(a,q)=1}\left|
 \pi(2x;q,a)-\pi(x;q,a)-\frac{\li(2x)-\li(x)}{\varphi(q)}\right|.
\]
The two endpoints follow from the usual maximal version with only a factor
two in the constant.  In particular $4x^{1/3}$ is below the level for
every fixed $R$ and all sufficiently large $x$.

We also use Shiu \cite[Theorem~1, pp.~162--163]{Shiu}.
If $F$ is nonnegative and multiplicative,
$F(p^a)\leq A_1^a$ for every prime power, and for each $\epsilon>0$
there is $A_2(\epsilon)$ with $F(n)\leq A_2(\epsilon)n^\epsilon$, then
for fixed $0<\alpha_s,\beta_s<1/2$,
\begin{equation}\label{eq:Shiu}
 \sum_{\substack{V-Y<n\leq V\\n\equiv a\ (q)}}F(n)
 \ll_{\alpha_s,\beta_s,F}\frac{Y}{\varphi(q)\log V}
 \exp\left\{\sum_{p\leq V,\ p\nmid q}\frac{F(p)}p\right\}
\end{equation}
uniformly for $q\geq2$, $(a,q)=1$, $q<Y^{1-\alpha_s}$ and
$V^{\beta_s}<Y\leq V$.  The constant includes the growth data of $F$.
In our applications $F(n)=a_0^{\omega(n)}$ for fixed $1<a_0<2$, or
$F(n)=n/\varphi(n)$.  Both satisfy these hypotheses.  For the first use
$F(n)\leq\tau(n)^{\log_2a_0}\ll_{a_0,\epsilon}n^\epsilon$;
for the second use the divisor expansion below.  Fix boxes
$(U,(1+\delta_0)U]$ with $\delta_0=1/20$, and take
$V=(1+\delta_0)U$, $Y=\delta_0 U$ and
$\alpha_s=\beta_s=1/4$.  For every $q\leq K$ and $U\geq H=K^{10}$,
both inequalities hold with room to spare once $K\geq K_0$.
The unrestricted $q=1$ upper estimates needed here follow elementarily
from $F=1*h$: for these two functions $h$ is nonnegative and supported
on squarefree integers, with $h(p)=a_0-1$ and $1/(p-1)$ respectively.
Indeed $\sum_{n\leq V}F(n)\leq V\sum_{d\leq V}h(d)/d$;
the Euler product gives $O_{a_0}(V(\log V)^{a_0-1})$ and $O(V)$.
Thus no endpoint convention about modulus one in Shiu is required.

For $k\geq1$, $(c,k)=1$, and $1\leq H<z$, put $\Lambda=\log(z/H)$ and set
\[
 \Sset(k,c;H,z)=\left\{(u,v):\begin{array}{l}
 H<u,v\leq z,\ (u,v)=1,\ (v,k)=1,\\
 u\equiv-cv\pmod k
 \end{array}\right\}.
\]
Here $(u,k)=1$ follows automatically.

\begin{lemma}[elementary lattice sums; {\cite[Lemma 3.1]{LL}}]\label{lem:lattice}
There exist absolute $K_0,C_1>0$ such that for every real $K\geq K_0$,
every integer $1\leq k\leq K$, every $(c,k)=1$, and every $z\geq H^2$
with $H=K^{10}$,
\begin{align}
 \sum_{\Sset(k,c;H,z)}\frac1{uv}
 &\geq\frac14\frac{\varphi(k)}{k^2}\Lambda^2,\label{eq:latlower}\\
 \sum_{\Sset(k,c;H,z)}\frac1{\varphi(u)\varphi(v)}
 &\leq C_1\frac{\varphi(k)}{k^2}\Lambda^2.\label{eq:latupper}
\end{align}
There is no polylogarithmic restriction on $K$.
\end{lemma}
\begin{proof}
We recall the elementary proof to make this uniformity explicit.  For
$H'\geq1$, $z'>H'$, $\Lambda'=\log(z'/H')$, and $(a,k)=1$, integral
comparison and M\"obius inversion give
\begin{align}
 \sum_{\substack{H'<u\leq z'\\u\equiv a\ (k)}}\frac1u
 &=\frac{\Lambda'}k+O(1/H'),\label{eq:apone}\\
 \sum_{\substack{H'<v\leq z'\\(v,k)=1}}\frac1v
 &=\frac{\varphi(k)}k\Lambda'+O(\tau(k)/H').\label{eq:aptwo}
\end{align}
These include empty progressions.  Write $T(H',z')$ for the double sum
without $(u,v)=1$.  For $\Lambda'\geq1$,
\[
 T(H',z')=\frac{\varphi(k)}{k^2}\Lambda'^2
              +O(\Lambda'\tau(k)/H').
\]
Inverting $(u,v)=1$, using $(d,k)=1$ for a common divisor, and separating
$d\leq H$ from $d>H$, gives
\[
 \sum_{\Sset}\frac1{uv}
 =\frac{\varphi(k)}{k^2}\Lambda^2
       \sum_{\substack{d\leq H\\(d,k)=1}}\frac{\mu(d)}{d^2}
 +O\left(\frac{\Lambda\tau(k)\log z}{H}+\frac{\Lambda^2}{H}\right).
\]
For the tail use $T(H/d,z/d)\leq(1+\log z)^2$ and
$\log z\leq2\Lambda$.  The main coefficient is
$\prod_{p\nmid k}(1-p^{-2})+O(H^{-1})\geq6/\pi^2-O(H^{-1})$.
Since $\tau(k)k^2/\varphi(k)\leq k\tau(k)^2\leq4K^2$, the relative
error is $O(K^2/H)=O(K^{-8})$.  This proves \eqref{eq:latlower}.

For the upper estimate drop $(u,v)=1$ and use
\[
 \frac1{\varphi(m)}=\frac1m\sum_{d\mid m}\frac{\mu^2(d)}{\varphi(d)}.
\]
Applying \eqref{eq:apone} for $d\leq H$ and the trivial harmonic bound
for $d>H$ shows, uniformly for reduced $a$,
\[
 \sum_{\substack{H<u\leq z\\u\equiv a\ (k)}}\frac1{\varphi(u)}
 \ll\frac\Lambda k+\frac{\log H}{H}+\frac\Lambda{\sqrt H}.
\]
Here $\sum_d\mu^2(d)/(d\varphi(d))<\infty$,
$\sum_{d\leq H}\mu^2(d)/\varphi(d)\ll\log H$, and
$\sum_{d>H}1/(d\varphi(d))\ll H^{-1/2}$.
For clarity the middle estimate follows by another use of the divisor
expansion and the convergent first sum; the tail follows from
$\varphi(d)\gg\sqrt d$.  These elementary bounds also prove convergence
of the first sum.  Similarly \eqref{eq:aptwo} gives
\[
 \sum_{\substack{H<v\leq z\\(v,k)=1}}\frac1{\varphi(v)}
 \ll\frac{\varphi(k)}k\Lambda+
        \frac{\tau(k)\log H}{H}+\frac\Lambda{\sqrt H}.
\]
Since $\tau(k)\leq2\sqrt K$, $\varphi(k)/k\geq1/(2\sqrt K)$ and
$\Lambda\geq10\log K$, the secondary terms are absorbed uniformly.
Multiply the last two estimates, applying the first with $a=-cv$, to
obtain \eqref{eq:latupper}.
\end{proof}

\begin{lemma}[multiplier harmonic sum; {\cite[Lemma 3.2]{LL}}]\label{lem:h}
For every real $K\geq5$,
\begin{equation}\label{eq:hfull}
 h(\K(K)):=\sum_{k\in\K(K)}\frac{\varphi(k)}{k^2}
       =\frac2{\pi^2}\log K+O(1),
\end{equation}
with absolute error constant.  For every prime $p\leq K$,
\begin{equation}\label{eq:hremove}
 \sum_{\substack{k\leq K\\p\mid k}}\frac{\varphi(k)}{k^2}
        \leq\frac{1+\log K}{p}.
\end{equation}
\end{lemma}
\begin{proof}
Expand $\varphi(k)/k=\sum_{d\mid k}\mu(d)/d$.  The progression
$k\equiv1\pmod4$ gives
\[
 h(\K(K))=\sum_{\substack{d\leq K\\d\ \mathrm{odd}}}\frac{\mu(d)}{d^2}
 \left(\frac14\log\frac Kd+O(1)\right).
\]
Now $\sum_{d\ \mathrm{odd}}\mu(d)/d^2=8/\pi^2$ and
$\sum_d(1+\log d)/d^2<\infty$.  The second assertion follows by
$\varphi(k)/k^2\leq1/k$ and writing $k=pm$.
\end{proof}

\begin{lemma}[low-$\omega$ mass in the full power range]\label{lem:omega}
There is a fixed absolute integer $D\geq1$ such that, for all $K\geq K_0$,
$H=K^{10}$, $z\geq H^2$ sufficiently large, $k\leq K$ and $(c,k)=1$,
\begin{equation}\label{eq:lowomega}
 \sum_{\substack{(u,v)\in\Sset(k,c;H,z)\\
                 \omega(uv)\leq D\log\log z}}\frac1{uv}
 \geq\frac18\frac{\varphi(k)}{k^2}\Lambda^2.
\end{equation}
The threshold and constants are absolute once $D$ is fixed.  Thus in
\eqref{eq:parameters}, uniformly over every $\J\subseteq\K(K)$ and every
$(c,L_{\J})=1$, summing \eqref{eq:lowomega} over $k\in\J$ gives a lower
bound $\gg(\log z)^2 h(\J)$, also with cutoff $D\log\log X$ when $z\leq X$.
Here $L_{\J}=\lcm_{k\in\J}k$ and $h(\J)=\sum_{k\in\J}\varphi(k)/k^2$.
\end{lemma}
\begin{proof}
Fix $a_0=3/2$.  Apply \eqref{eq:Shiu} in the boxes described above to
$F(n)=a_0^{\omega(n)}$.  For each reduced $v$ the $u$-congruence is
reduced; summing the second variable over all reduced classes gives,
after division by $UV$, a bound
\[
 C\frac1{\varphi(k)}(\log U\log V)^{a_0-1}
                    \exp\{-2a_0\sum_{p\mid k}p^{-1}\}.
\]
The ratio of its local coefficient to $\varphi(k)/k^2$ has logarithm
\[
 \sum_{p\mid k}\{-2\log(1-1/p)-2a_0/p\}
 \leq C_{a_0}\sum_p p^{-2};
\]
the linear terms are nonpositive.  All primes dividing $k$ are below the
box endpoints, so none is lost from this calculation.  Drop $(u,v)=1$
for this upper estimate, cover the interval by full boxes (enlarging the
last), and sum $O(\Lambda^2)$ box pairs.  It follows that
\begin{equation}\label{eq:omegamoment}
 \sum_{\Sset}\frac{a_0^{\omega(u)+\omega(v)}}{uv}
 \ll_{a_0}\frac{\varphi(k)}{k^2}\Lambda^2
                       (\log z)^{2(a_0-1)}.
\end{equation}
On the coprime pairs $\omega(uv)=\omega(u)+\omega(v)$.  The omitted tail
is consequently at most the right side of \eqref{eq:omegamoment}
multiplied by $(\log z)^{-D\log a_0}$.  Fix any integer $D$ with
$D\log(3/2)>1+2$ and enlarge the threshold until this tail is at most
one eighth of $\varphi(k)\Lambda^2/k^2$.
Lemma~\ref{lem:lattice} then proves the claim.  No application of Shiu
has growing function parameters or an unreduced residue, and no step
requires $K$ to be a power of a logarithm.
\end{proof}

\section{Congestion pruning and cubic prime-slice mass}\label{sec:harvest}
For $(c,L_{\J})=1$, define
\[
 r_{\J}(u,v;c)=\#\{k\in\J:k\mid u+cv\}.
\]
If $(u,v)=1$ and $k\mid u+cv$, reducedness of $c$ modulo $k$ implies
$(uv,k)=1$.  Thus this incidence counts precisely the multiplier choices
in the lattice, before the prime progression is imposed.

\begin{lemma}[second incidence moment; notes Lemma 34.7]\label{lem:congestion}
There is an absolute constant $C_2$ such that for every $K\geq K_0$,
$H=K^{10}$, $z\geq H^2$, every $\J\subseteq\K(K)$ and every
$(c,L_{\J})=1$,
\begin{equation}\label{eq:second}
 \sum_{\substack{H<u,v\leq z\\(u,v)=1}}\frac{r_{\J}(u,v;c)^2}{uv}
 \leq C_2(\log z)^2(1+\log K)^3.
\end{equation}
\end{lemma}
\begin{proof}
Expand the square.  A pair $k,k'$ requires the single congruence
$d=[k,k']\mid u+cv$, with $(c,d)=1$ and $d\leq K^2$.
The same calculation \eqref{eq:apone}--\eqref{eq:aptwo}, now at modulus
$d$ and without the coprimality restriction on the pair, gives
\begin{equation}\label{eq:lcmdensity}
 \sum_{\substack{H<u,v\leq z\\(u,v)=1,\ d\mid u+cv}}\frac1{uv}
 \ll\frac{\varphi(d)}{d^2}(\log z)^2+
                       \frac{\tau(d)(\log z)^2}{H}.
\end{equation}
Indeed the original coprime pairs have $(v,d)=1$, so dropping $(u,v)=1$
while retaining $(v,d)=1$ is an upper bound.  The direct harmonic
calculation gives even $O(\Lambda\tau(d)/H)$ here; the weaker displayed
error is the box-boundary estimate in notes (34.15).  This does not
misapply Lemma~\ref{lem:lattice} with floor $(K^2)^{10}$.

For the main coefficients,
\[
 \sum_{k,k'\leq K}\frac{\varphi([k,k'])}{[k,k']^2}
 \leq\sum_{k,k'\leq K}\frac{(k,k')}{kk'}
 \leq\sum_{d\leq K}\frac1d(1+\log(K/d))^2
 \ll(1+\log K)^3.
\]
Here writing $k=da,k'=db$, $(a,b)=1$, and then dropping this last
condition proves the middle bound.  The total errors in
\eqref{eq:lcmdensity} are $O(K^3(\log z)^2/H)$, because
$\tau([k,k'])\leq2K$ and there are at most $K^2$ pairs.
Since $H=K^{10}$ these errors are absorbed in \eqref{eq:second}.
\end{proof}

We state the unused hypothesis to distinguish precisely what pruning proves.
The literal $H_{k\mathrm{BV}}(\kappa)$ of notes (34.2) asserts the existence
of a function $\rho_\kappa(X)\to0$ such that, for all sufficiently large
$X$, simultaneously for $X^{1/2}\leq x\leq X$, $K_0\leq K\leq X^\kappa$,
every $1\in\J\subseteq\K(K)$ and every reduced $c\pmod{24L_{\J}}$,
\begin{equation}\label{eq:HkBV}
 \sum_{k\in\J}\sum_{\substack{(u,v)\in\Sset(k,c;H,z)\\
                         \omega(uv)\leq D\log\log X}}
 |E(x;4uv,-k^{-1})|
 \leq\rho_\kappa(X)x\log x\,h(\J),
\end{equation}
where $H=K^{10}$, $z=x^{1/6}$ and $E$ is the error defining $E_x^*$.
Every triple retains its multiplicity.  The reducedness modulo $24L_{\J}$
can be replaced in the harvesting proof by reducedness modulo $L_{\J}$:
since $L_{\J}$ is odd, every such class lifts to a reduced class modulo
$24L_{\J}$, without changing any multiplier congruence.
This hypothesis is \emph{not}
assumed or proved.  Plain BV would multiply \eqref{eq:BV} by as much as
$K(\log X)^{D\log2}$, which cannot be absorbed for $K=X^\kappa$ by a
fixed logarithmic saving.

\begin{theorem}[pruned prime slice; notes Theorem 34.8]\label{thm:harvest}
Fix $0<\kappa<1/240$ and $D$ from Lemma~\ref{lem:omega}.  There are
$a_h,A_h>0$ and $X_h=X_h(\kappa,D)$ such that for every $X\geq X_h$,
every $K_0\leq K\leq X^\kappa$, every $1\in\J\subseteq\K(K)$, and
 every $(c,L_{\J})=1$, the following holds.  Within each canonical block
keep only triples with
\begin{equation}\label{eq:pruning}
 \begin{gathered}
 (u,v)\in\Sset(k,c;H,z_j),\quad k\in\J,\quad
 \omega(uv)\leq D\log\log X,\\
 r_{\J}(u,v;c)\leq T_X:=t^4,\quad4uv\mid k\ell+1.
 \end{gathered}
\end{equation}
Let $f^{\rm good}_{c;\J}(\ell)$ count their distinct classes modulo $\ell$.
Then uniformly in all these data,
\begin{equation}\label{eq:harvest}
 a_ht^2h(\J)\leq\sum_{X^{1/2}<\ell\leq X}
                  \frac{f^{\rm good}_{c;\J}(\ell)}\ell
                    \leq A_ht^2h(\J).
\end{equation}
Moreover \eqref{eq:HkBV} restricted by $r_{\J}\leq T_X$ is
$o(x\log x\,h(\J))$ uniformly, without assuming $H_{k\mathrm{BV}}$.
\end{theorem}
\begin{proof}
In a full block put $z=x^{1/6}$.  By \eqref{eq:floor} and
Lemma~\ref{lem:omega} the low-$\omega$ harmonic triple mass is
$\gg(\log z)^2h(\J)$.  The mass removed by the congestion condition is
at most
\begin{equation}\label{eq:prunedmass}
 \frac1{t^4}\sum_{u,v}\frac{r_{\J}(u,v;c)^2}{uv}
 \ll\frac{(\log z)^2(1+\log K)^3}{t^4}
 =o((\log z)^2h(\J)).
\end{equation}
Uniformity for sparse subfamilies uses $h(\J)\geq1$ from $1\in\J$;
indeed the relative error is $O_\kappa(1/t)$ even in that worst case.
This proves the negligible loss of \emph{harmonic} mass.  To obtain the
same assertion for the actual prime main-term weight, use the elementary
Mertens consequence $n/\varphi(n)\ll\log\log(3n)$: split its prime product
at $\log(3n)$; the smaller primes cost $O(\log\log(3n))$ and the larger
ones a bounded factor.  Since $4uv\leq4X^{1/3}$, the removed
$1/\varphi(4uv)$ mass is at most $O(\log t)$ times the removed harmonic
mass.  Relative to the original main-term mass, which is
$\gg(\log z)^2h(\J)$, its loss is $O_\kappa(\log t/t)=o(1)$.
Thus both weighted masses are controlled; no assertion is made that the
discarded \emph{actual prime count} is small.

Each retained triple asks that
\begin{equation}\label{eq:primeAP}
 \ell\equiv-k^{-1}\pmod{4uv},\qquad q=4uv\leq4x^{1/3}.
\end{equation}
The residue is reduced and already is $3\pmod4$.  For a fixed $q=4uv$,
coprimality allows at most $2^{\omega(uv)}$ ordered allocations of its
prime powers to $u,v$.  Each allocation has at most $T_X$ admissible
multipliers.  Hence its entire retained multiplicity is
\begin{equation}\label{eq:Wgood}
 W_{\rm good}(q)\leq t^{4+D\log2}.
\end{equation}
Choose once and for all $R>4+D\log2+10$ in \eqref{eq:BV}.
Since $\log x\asymp t$ and $q\leq4x^{1/3}$ lies below the BV level,
grouping by $q$ bounds the whole retained triple error by
$O_D(x(\log x)^{-10})$.  All constants and thresholds here are independent
of $c,\J,K,x$ in their specified ranges.

For the main term use $\li(2x)-\li(x)\geq x/\log(2x)$ and
$1/\varphi(4uv)\geq1/(4uv)$.  Equation~\eqref{eq:prunedmass} therefore
gives at least $a x\log x\,h(\J)$ retained incidences.  Lemma~\ref{lem:CRT}
proves that these are genuinely distinct classes for each $\ell$; in
particular repeated occurrences of a prime are not an overcount of a
single class.  Dividing by $2x$ gives $\gg\log x\,h(\J)$ reciprocal mass
in each full block.

For the upper bound drop both prunings.  Brun--Titchmarsh and
$\varphi(4uv)\geq2\varphi(u)\varphi(v)$ for coprime $u,v$ give
\[
 \sum_{x<\ell\leq2x} f_{c;\J}(\ell)
 \ll\frac{x}{\log x}\sum_{k\in\J}\sum_{\Sset(k,c;H,z)}
                            \frac1{\varphi(u)\varphi(v)}
 \ll x\log x\,h(\J)
\]
by Lemma~\ref{lem:lattice}.  Divide by $x$ and sum over the blocks.
There are $\asymp t$ full blocks for the lower bound; enlarging the last
block gives the upper bound.  This proves all assertions.
\end{proof}

\begin{corollary}[uniform fibre masses]\label{cor:fibremass}
For every $X\geq X_h$ with $K=\lfloor X^\kappa\rfloor$ and every residue
$c\pmod{L_K}$, put
\[
 \begin{gathered}
 \J_c=\{k\in\K(K):(k,c)=1\},\qquad\mu_c=\sum_\ell f_c(\ell)/\ell,\\
 f_c(\ell)=\#\{A\in\A_X:\ell_A=\ell,\ k_A\mid u_A+cv_A\}.
 \end{gathered}
\]
There are constants $a_h,C_u>0$, depending at most on $\kappa,D$, such that
\begin{equation}\label{eq:fibremass}
 a_ht^2h(\J_c)\leq\mu_c\leq A_ht^2h(\J_c)\leq C_ut^3
\end{equation}
uniformly in \emph{all} $c$, not only reduced residues modulo $L_K$.
For reduced $c$ the mass is $\asymp_\kappa t^3$.
\end{corollary}
\begin{proof}
An active atom necessarily has $(k,c)=1$, since its residue modulo $k$ is
a unit.  Thus active atoms have multipliers in $\J_c$, which contains 1,
and $(c,L_{\J_c})=1$.  The lower bound of Theorem~\ref{thm:harvest}
applies to this very subfamily, and its canonical good triples are contained
in $\A_X$.  The unpruned BT upper bound in that proof applies with the
same $\J_c$.  Finally use Lemma~\ref{lem:h} and $\log K\leq\kappa t$.
For reduced $c$, $\J_c=\K(K)$ and $\log K\geq\kappa t/2$ once $X$ is
large.  This checks explicitly the data-dependent-subfamily quantifiers.
\end{proof}


\section{Unpruned supply by Cauchy--Schwarz (simplification)}\label{sec:unpruned}
This section, added after the independent re-derivation recorded in
\cite[\S76]{Notes}, gives Theorem~\ref{thm:harvest} for the
\emph{full} triple family, with no cutoff on $\omega(uv)$ and no
congestion pruning.  With it, Lemma~\ref{lem:omega}, Lemma~\ref{lem:congestion}
as a pruning device, the constant $D$, and the cutoff $T_X$ are no
longer used anywhere in the proof of Theorem~\ref{thm:main}
(Corollary~\ref{cor:noD}).  Section~\ref{sec:harvest} is retained as
the reviewed original route.

\begin{lemma}[weighted second incidence moment]\label{lem:wsecond}
There is an absolute constant $C_3$ such that for every $K\geq K_0$,
$H=K^{10}$, $z\geq H^2$, every $\J\subseteq\K(K)$ and every
$(c,L_{\J})=1$,
\begin{equation}\label{eq:wsecond}
 \sum_{\substack{H<u,v\leq z\\(u,v)=1}}
 \frac{2^{\omega(uv)}\,r_{\J}(u,v;c)^2}{\varphi(4uv)}
 \leq C_3(\log z)^4(1+\log K)^3.
\end{equation}
\end{lemma}
\begin{proof}
For coprime $u,v$ one has $\varphi(4uv)\geq2\varphi(u)\varphi(v)$: if
$uv$ is odd this is an equality, and if $u=2^am$ with $a\geq1$, $m$ odd,
then $\varphi(4u)=2^{a+1}\varphi(m)=4\varphi(u)$.  Expand the square as
in Lemma~\ref{lem:congestion}: a pair $k,k'\in\J$ imposes the single
congruence $d=[k,k']\mid u+cv$ with $d\leq K^2$ and $(c,d)=1$.  For
coprime $u,v$ with $d\mid u+cv$ one has $(v,d)=1$, so dropping $(u,v)=1$
while keeping $(v,d)=1$ is an upper bound, and the residue
$u\equiv-cv\pmod d$ is reduced.  Thus each pair contributes at most
\[
 \tfrac12\sum_{\substack{H<v\leq z\\(v,d)=1}}\frac{2^{\omega(v)}}{\varphi(v)}
 \sum_{\substack{H<u\leq z\\u\equiv-cv\ (d)}}\frac{2^{\omega(u)}}{\varphi(u)}.
\]
Cover $(H,z]$ by consecutive boxes $(U,(1+\delta_0)U]$ starting at
$U=H$, enlarging the last box when necessary; all endpoints are then at
most $(21/20)z$, which affects no logarithmic bound.  For $d\geq2$ apply
\eqref{eq:Shiu} to $F(n)=2^{\omega(n)}n/\varphi(n)$, which is
multiplicative with $F(p^a)=2p/(p-1)\leq4\leq4^a$ and
$F(n)\leq4^{\omega(n)}\leq\tau(n)^2\ll_\epsilon n^\epsilon$, in these full
boxes, with modulus $d\leq K^2\leq U^{1/5}$ and a reduced residue.  For
$d=1$ (the pair $k=k'=1$, always present since $1\in\J$), the
nonnegative convolution $F=1*h$ below gives $\sum_{n\leq V}F(n)\ll V\log V$,
hence each box contributes at most $U^{-1}\sum_{n\leq(21/20)U}F(n)\ll\log U$,
which is the same estimate since $\varphi(1)/1^2=1$.  Since
$\sum_{p\leq V}F(p)/p=2\log\log V+O(1)$ and
$\exp\{-\sum_{p\mid d}2/(p-1)\}\asymp(\varphi(d)/d)^2$ (the logarithm
of the ratio is $\sum_{p\mid d}O(p^{-2})$), the inner box sum of
$2^{\omega(u)}/\varphi(u)\leq F(u)/U$ is $\ll(\log U)\varphi(d)/d^2$.
There are $O(\log z)$ boxes, so the inner sum is
$\ll(\log z)^2\varphi(d)/d^2$, uniformly in the residue.  For the outer
sum drop $(v,d)=1$; $F=1*h$ with $h\geq0$ supported on squarefree
integers, $h(p)=(p+1)/(p-1)$ (indeed $h(p^a)=F(p^a)-F(p^{a-1})=0$ for
$a\geq2$), gives $\sum_{n\leq V}F(n)\ll V\log V$ and,
by partial summation, $\sum_{H<v\leq z}2^{\omega(v)}/\varphi(v)\ll(\log z)^2$.
Hence each pair contributes $\ll(\log z)^4\varphi(d)/d^2$, and
$\sum_{k,k'\leq K}\varphi([k,k'])/[k,k']^2\ll(1+\log K)^3$ exactly as in
Lemma~\ref{lem:congestion}.
\end{proof}

For $1\in\J\subseteq\K(K)$ and $(c,L_{\J})=1$ let $f^{\rm all}_{c;\J}(\ell)$
count the triples $(k,u,v)$ with $k\in\J$, $(u,v)\in\Sset(k,c;H,z_j)$
for the block $I_j\ni\ell$, and $4uv\mid k\ell+1$.  By
Lemma~\ref{lem:CRT}, whose distinctness argument uses coprimality,
$u,v\leq z_j$ with $z_j^2<\ell$, $H<u,v$ with $4H^2>K$, and
$4uv\mid k\ell+1$, but no $\omega$ cutoff, these triples have
pairwise distinct residues modulo $\ell$, so $f^{\rm all}$ also counts
distinct classes.

\begin{theorem}[unpruned prime slice]\label{thm:unpruned}
Fix $0<\kappa<1/240$.  There are $a_h,A_h>0$ and $X_h=X_h(\kappa)$ such
that for every $X\geq X_h$, every $K_0\leq K\leq X^\kappa$, every
$1\in\J\subseteq\K(K)$ and every $(c,L_{\J})=1$,
\begin{equation}\label{eq:unpruned}
 a_ht^2h(\J)\leq\sum_{X^{1/2}<\ell\leq X}\frac{f^{\rm all}_{c;\J}(\ell)}{\ell}
 \leq A_ht^2h(\J).
\end{equation}
\end{theorem}
\begin{proof}
The upper bound is the last paragraph of the proof of
Theorem~\ref{thm:harvest}, which used no pruning.  For the lower bound
fix a full block $(x,2x]$ with $z=x^{1/6}$, so that $4uv\leq4x^{1/3}$
and $\Lambda\geq\frac12\log z\asymp t$ by \eqref{eq:floor}.  Then
\[
 \sum_{x<\ell\leq2x}f^{\rm all}_{c;\J}(\ell)
 =\sum_{k\in\J}\sum_{(u,v)\in\Sset(k,c;H,z)}
 \bigl[\pi(2x;4uv,-k^{-1})-\pi(x;4uv,-k^{-1})\bigr],
\]
the residue being reduced since $(k,4uv)=1$.  Write each bracket as
$(\li(2x)-\li(x))/\varphi(4uv)+E(x;4uv,-k^{-1})$.  By \eqref{eq:latlower}
the main term is at least
\[
 \frac{x}{\log(2x)}\sum_{k\in\J}\sum_{\Sset(k,c;H,z)}\frac1{4uv}
 \geq\frac{x}{16\log(2x)}\Lambda^2h(\J)\gg x\,t\,h(\J).
\]
For the error, group the triples by $q=4uv$ and put
$W_c(q)=\sum_{(u,v):\,4uv=q}r_{\J}(u,v;c)$, the sum over coprime
$H<u,v\leq z$ with $uv=q/4$; there are at most $2^{\omega(q/4)}$ such
ordered pairs, so by Cauchy--Schwarz
$W_c(q)^2\leq2^{\omega(q/4)}\sum_{4uv=q}r_{\J}(u,v;c)^2$.  A second
Cauchy--Schwarz gives
\[
 \Bigl|\sum_{\rm triples}E\Bigr|\leq\sum_qW_c(q)E^*_x(q)
 \leq\Bigl(\sum_qW_c(q)^2E^*_x(q)\Bigr)^{1/2}
 \Bigl(\sum_{q\leq4x^{1/3}}E^*_x(q)\Bigr)^{1/2}.
\]
For $q\leq4x^{1/3}$, \eqref{eq:BT} and $\li(2x)-\li(x)\leq x/\log x$ give
$E^*_x(q)\ll x/(\varphi(q)\log x)$, so by Lemma~\ref{lem:wsecond}
\[
 \begin{aligned}
 \sum_qW_c(q)^2E^*_x(q)
 &\ll\frac{x}{\log x}
 \sum_{\substack{H<u,v\leq z\\(u,v)=1}}
 \frac{2^{\omega(uv)}r_{\J}(u,v;c)^2}{\varphi(4uv)}\\
 &\ll\frac{x}{\log x}(\log z)^4(1+\log K)^3\ll x\,t^6,
 \end{aligned}
\]
using $\log z\asymp\log x\asymp t$ and $\log K\leq\kappa t$.  With the
fixed choice $R=26$ in \eqref{eq:BV},
$\sum_{q\leq4x^{1/3}}E^*_x(q)\leq C_{26}x(\log x)^{-26}$ for
$x\geq x_{26}$.  The total error is therefore $\ll x\,t^{-10}$, which is
$o(x\,t\,h(\J))$ because $h(\J)\geq1$ from $1\in\J$.  Divide by $2x$ and
sum over the $\asymp t$ full blocks.  All constants depend only on
$\kappa$ and on the absolute Shiu, Brun--Titchmarsh and
Bombieri--Vinogradov constants; none depends on $c$, $\J$, $K$ or $x$.
\end{proof}

\begin{corollary}[the theorem without $D$, pruning, or
Lemma~\ref{lem:congestion}]\label{cor:noD}
Let $\A'_X$ be the atom family \eqref{eq:atomdef} with the condition
$\omega(uv)\leq D\log\log X$ deleted.  Then Lemma~\ref{lem:CRT},
Corollary~\ref{cor:fibremass}, Theorem~\ref{thm:moments}
(conditional-independence proof), Lemma~\ref{lem:void},
Theorem~\ref{thm:assembly} and Section~\ref{sec:all} hold with $\A_X$
replaced by $\A'_X$, Theorem~\ref{thm:harvest} replaced by
Theorem~\ref{thm:unpruned}, the references to a pruned lower estimate
replaced by the unpruned supply estimate, and $D$ deleted from all
parameter dependencies; and Theorem~\ref{thm:main} follows with the
same order for the surviving parameters and without any use of
Lemma~\ref{lem:omega}, Lemma~\ref{lem:congestion}, the lower $W_k$
calculation of Lemma~\ref{lem:profile}, the ordered-atom replay, $D$,
or $T_X$.
\end{corollary}
\begin{proof}
In the original route, the cutoff and the congestion threshold are
used to establish low-$\omega$ mass (Lemma~\ref{lem:omega}), to control
the retained main-term mass \eqref{eq:prunedmass}, and to bound the
Bombieri--Vinogradov multiplicity \eqref{eq:Wgood}; the optional lower
$W_k$ calculation in Lemma~\ref{lem:profile} also uses the cutoff.
Theorem~\ref{thm:unpruned} replaces the entire pruned-supply argument.
The proof of the main theorem then uses only the conditional-independence
proof of Theorem~\ref{thm:moments}, not Lemma~\ref{lem:profile} or the
ordered-atom replay.  Lemma~\ref{lem:CRT} does not use the cutoff.  Corollary~\ref{cor:fibremass}
uses the lower bound of the supply theorem for the subfamily
$\J_c\ni1$ with $(c,L_{\J_c})=1$, and the unpruned upper bound; both
are \eqref{eq:unpruned}.  The conditional-independence proof of
Theorem~\ref{thm:moments} needs only Corollary~\ref{cor:fibremass};
Lemma~\ref{lem:void} needs only the fibre lower bound and the selector
computation; Theorem~\ref{thm:assembly} needs the moment bound, the
void bound, and the crude inventory \eqref{eq:inventory}, which never
used the cutoff.  The transfer to all denominators is untouched.
\end{proof}

\begin{remark}
The device is the one used for the fixed-multiplier theorem in
\cite[Section~4]{LL}: a multiplicity as large as $X^\kappa$ at one
modulus $q$ is harmless when the error is controlled in mean square,
because it is paid for by its own weight $1/\varphi(q)$ in
$\sum_qW_c(q)^2/\varphi(q)$.  In the power range the weight
$2^{\omega(uv)}$ makes the allocation count explicit, and Shiu's
theorem for $2^{\omega(n)}n/\varphi(n)$ costs two extra powers of
$\log z$, which the fixed logarithmic saving of Bombieri--Vinogradov
absorbs.  The unused hypothesis \eqref{eq:HkBV} is now irrelevant to the
proof rather than merely unassumed.  The exponent, the cubic mass per
fibre, the Bonferroni degree and the ledger are unchanged.
\end{remark}

\section{Residue profiles and factorial moments}\label{sec:moments}
This section retains the residue-resolved estimate and the local induction
of notes \S\S39.2--39.3, including the referee-sensitive factors.  We also
give the conditional-independence formulation of the same moment bound;
Corollary~\ref{cor:fibremass} makes its uniform upper hypothesis explicit.

For $k\in\K(K)$, $g\mid k$, and a reduced residue $a\pmod g$, define
\[
 W_{k,a}(g)=\sum_{\substack{A\in\A_X,\ k_A=k\\-u_Av_A^{-1}\equiv a\ (g)}}
                    \frac1{k\ell_A},\qquad
 W_k=\sum_{\substack{A\in\A_X\\k_A=k}}\frac1{k\ell_A}.
\]
The sole class modulo $g=1$ is understood as reduced.

\begin{lemma}[residue-resolved upper bound; notes (39.11)]\label{lem:profile}
For fixed $\kappa,D$ there are $a_p,A_p>0$ and $X_p=X_p(\kappa,D)$ such
that for every $X\geq X_p$, every $k\in\K(K)$, every $g\mid k$ and every
reduced $a\pmod g$,
\begin{equation}\label{eq:profiles}
 W_{k,a}(g)\leq\frac{A_pt^2}{\varphi(g)}\frac{\varphi(k)^2}{k^3},
 \qquad
 a_pt^2\frac{\varphi(k)^2}{k^3}\leq W_k
       \leq A_pt^2\frac{\varphi(k)^2}{k^3}.
\end{equation}
In particular $\mu_X:=\sum_A1/(k_A\ell_A)\asymp_\kappa t^3$.
\end{lemma}
\begin{proof}
Fix a block and put $q_0=\lcm(g,\rad k)\leq k$.  In boxes
$u\asymp U$, $v\asymp V$, the conditions $(uv,k)=1$ and
$-uv^{-1}\equiv a\pmod g$ permit exactly
$\varphi(q_0)^2/\varphi(g)$ ordered pairs of unit classes modulo $q_0$.
Indeed reduction of units modulo $q_0$ onto units modulo $g$ is
surjective, with all fibres of size $\varphi(q_0)/\varphi(g)$.
Apply Shiu to $F(n)=n/\varphi(n)$ in each class.  Its prime value is
$p/(p-1)$, whence Mertens' estimate gives
\[
 \frac1{\log U}\exp\left\{\sum_{p\leq(1+\delta_0)U,\ p\nmid q_0}
                              \frac1{p-1}\right\}
 \ll\exp\left\{-\sum_{p\mid k}\frac1{p-1}\right\}
 \asymp\frac{\varphi(k)}k.
\]
The last comparability is uniform: the logarithm of its ratio is a sum
of $O(p^{-2})$.  Shiu applies since $q_0\leq K\leq U^{1/10},V^{1/10}$,
with the interval conditions already checked in Section~\ref{sec:lattice}.
Drop only $(u,v)=1$ for an upper bound.  Multiplying the two one-class
estimates by the number of compatible pairs and dividing by $UV$ gives
\[
 \sum_{\rm box}\frac1{\varphi(u)\varphi(v)}
 \ll\frac1{\varphi(g)}\left(\frac{\varphi(k)}k\right)^2.
\]
Summing the $O((\log z)^2)$ box pairs proves the full residue estimate
\begin{equation}\label{eq:residuebox}
 \sum_{\substack{H<u,v\leq z,\ (uv,k)=1\\-uv^{-1}\equiv a\ (g)}}
 \frac1{\varphi(u)\varphi(v)}
 \ll\frac{(\log z)^2}{\varphi(g)}\frac{\varphi(k)^2}{k^2}.
\end{equation}
No divisor factor or $k^\epsilon$ is suppressed.  When applying BT we
restrict back to coprime pairs, so that
$1/\varphi(4uv)\leq1/(2\varphi(u)\varphi(v))$ is valid.
Equation~\eqref{eq:BT} implies
\[
 \sum_{\substack{x<\ell\leq2x\\\ell\equiv-k^{-1}\ (4uv)}}\frac1\ell
 \ll\frac1{\varphi(4uv)\log x}.
\]
Combine this with \eqref{eq:residuebox}, divide by $k$ and sum $O(t)$
blocks to get the first inequality of \eqref{eq:profiles}.

For completeness the lower $W_k$ bound does not assume a polynomially
small BV main term can survive a fixed logarithmic error.  For a fixed $k$,
sum Lemma~\ref{lem:omega} over the $\varphi(k)$ residues $c\pmod k$.
Each reduced coprime pair occurs in exactly one of them, giving low-$\omega$
harmonic mass $\gg(\log z)^2\varphi(k)^2/k^2$ without a ratio condition.
For this fixed $k$, a modulus has at most $t^{D\log2}$ ordered allocations;
there is no factor $K$.  BV with fixed $R>D\log2+10$ makes its error
$O_D(x(\log x)^{-10})$.
The elementary Mertens product implies
$\varphi(k)/k\geq\prod_{p\leq K}(1-1/p)\gg1/\log K$ for $K\geq3$.
Thus the block main term
$\gg x\log x\,(\varphi(k)/k)^2$ is uniformly $\gg_\kappa x/t$,
and dominates that error.  Division by $2xk$ and summation over full
blocks prove the lower bound.

Finally, writing $b(k)=(\varphi(k)/k)^2$, Cauchy--Schwarz and
Lemma~\ref{lem:h} give
\[
 \sum_{k\in\K(K)}\frac{b(k)}k
 \geq\frac{h(\K(K))^2}{\sum_{k\in\K(K)}1/k}\gg\log K;
\]
the reverse inequality follows from $b(k)\leq1$.  This proves the mass
assertion.  Notice that $\mu_X$ is an \emph{unconditional cylinder mass},
whereas $\mu_c$ in \eqref{eq:fibremass} is a conditional prime-coordinate
mass.  They are not interchangeable term by term.
\end{proof}

At an odd prime $p$ the nonconstant part of the proxy Euler factor for
$b(k)/k$ is
\[
 \sum_{e\geq1}\frac{(1-1/p)^2}{p^e}=\frac{p-1}{p^2}.
\]
Its share of the full local factor is $(p-1)/(p^2+p-1)$.
This is an exact \emph{Euler-factor} identity, not an exact finite-$K$
share of $W_k$, for which only comparability has been proved.

Define, for $y\geq2$,
\[
 q_y(k)=\prod_{p\mid k,\ p\leq y}\frac p{p-1},\qquad
 b_y(g)=\prod_{p^e\Vert g,\ p\leq y}\varphi(p^e)
                    \prod_{p^e\Vert g,\ p>y}p^e.
\]
\begin{lemma}[uniform multiplier bound; notes Lemma 39.3]\label{lem:euler}
For every real $y\geq2$, integer $R\geq1$, and real $K\geq3$,
\begin{equation}\label{eq:euler}
 \sum_{k\in\K(K)}\frac{\varphi(k)^2}{k^3}q_y(k)
            \prod_{p\mid(k,R),\ p>y}\frac p{p-1}\ll\log K,
\end{equation}
with an absolute constant independent of $R,y$.
\end{lemma}
\begin{proof}
Enlarge to all odd $k$ and to their Euler product with $p\leq K$.
An unmodified nonconstant local mass is $(p-1)/p^2$; either permitted
modification changes it to $1/p$.  The ratio of the modified full factor
to the original is $1+O(p^{-2})$, with a convergent product.  The original
product is bounded by $\prod_{p\leq K}(1-1/p)^{-1}\ll\log K$.
All terms are nonnegative, so the product enlargement is legitimate.
\end{proof}

\begin{theorem}[all upper factorial moments; notes Theorem 39.4]\label{thm:moments}
For fixed $\kappa,D$ there are $C_F,X_F>0$ such that, for every
$X\geq X_F$, every real $2\leq y<X^{1/2}$, and \emph{every} integer
$m\geq1$,
\begin{equation}\label{eq:moments}
 \E((H_X)_m\mid S_y=1)\leq(C_Ft^3)^m.
\end{equation}
Here $S_y=\ind_{(n,P_y)=1}$ and $(h)_m=h(h-1)\cdots(h-m+1)$.
In particular the base $C_F$ is independent of the order $m$.
\end{theorem}
\begin{proof}[Conditional-independence proof]
Reveal $c\pmod{L_K}$, retaining the condition $S_y=1$.
Since $y<X^{1/2}$, this does not condition any large-prime coordinate.
By Lemma~\ref{lem:CRT}, in this fibre $H_X$ is a sum of independent
Bernoulli variables $I_\ell$ with means $f_c(\ell)/\ell$; at most one
atom at each $\ell$ can be hit.  Expanding the falling factorial gives
\[
 \E((H_X)_m\mid c,S_y=1)
 =\sum_{\ell_1,\ldots,\ell_m\ \mathrm{distinct}}
                 \prod_{i=1}^m\frac{f_c(\ell_i)}{\ell_i}
 \leq\mu_c^m\leq(C_ut^3)^m.
\]
Corollary~\ref{cor:fibremass} applies to every $c$ admitted by the selector.
Averaging over these fibres proves \eqref{eq:moments} with $C_F=C_u$.
This conditioning is solely in the CRT probability space and creates
no finite-interval coefficient cost.
\end{proof}

For transparency we also replay the notes' ordered-atom proof, so the
third referee checkpoint is not bypassed by the formulation just given.
Given previous compatible atoms with multiplier lcm $R$, a new atom with
multiplier $k$, $g=(k,R)$, and a new $\ell$, has relative probability
\begin{equation}\label{eq:localrelative}
 \frac{q_y(k)}{k\ell}b_y(g)
\end{equation}
if its unit residue agrees modulo $g$, and zero otherwise.  Here is the
prime-power check.  If $p^e\Vert k$ and $p^f\Vert R$ (allow $f=0$),
the new local probability, subject to agreement, is
\[
\begin{array}{c|c}
 f\geq e&1\\
 0<f<e&p^{-(e-f)}\\
 f=0,\ p\leq y&1/\varphi(p^e)\\
 f=0,\ p>y&p^{-e}.
\end{array}
\]
For $p\leq y$ the factor in \eqref{eq:localrelative} is
$[p/(p-1)]p^{-e}\varphi(p^{\min(e,f)})$ when $f>0$;
this is precisely $p^{\min(e,f)-e}$.  The other cases follow immediately.
Unequal prime-power exponents are thus included, not replaced by radicals.

Compatibility fixes one unit residue modulo $g$.  Apply
Lemma~\ref{lem:profile} to its atom sum, multiply by $q_y(k)b_y(g)$,
and use the exact cancellation
\begin{equation}\label{eq:cancel}
 \frac{b_y(g)}{\varphi(g)}
       =\prod_{p\mid g,\ p>y}\frac p{p-1}.
\end{equation}
The sum of the relative probabilities of all permissible new atoms is
at most
\[
 A_pt^2\sum_{k\in\K(K)}\frac{\varphi(k)^2}{k^3}q_y(k)
               \prod_{p\mid(k,R),\ p>y}\frac p{p-1}
 \ll t^2\log K\ll_\kappa t^3
\]
by Lemma~\ref{lem:euler}.  Excluding old $\ell$ only decreases this sum.
Iteration over ordered compatible distinct atoms proves
\eqref{eq:moments} again, with another fixed base constant.
Discarding consistency would instead charge shared primes by their full
lcm-collapse factors.  That estimate is too large; a naive Janson
argument is neither used nor justified here.

\section{An exponentially small conditioned void}\label{sec:void}
The all-integer CRT distribution contains nonreduced multiplier fibres,
even though exceptional primes larger than $K$ are reduced.  One cannot
replace this distribution by a uniform distribution on reduced residues
without paying for that replacement.  We use the exact small-prime selector
\begin{equation}\label{eq:selector}
 y=Bt^3,\qquad S_y(n)=\ind_{(n,P_y)=1}
       =\sum_{d\mid P_y}\mu(d)\ind_{d\mid n}.
\end{equation}
For any fixed $B$, eventually $2\leq y<K<X^{1/2}$.  Its expansion has
$2^{\pi(y)}$ terms, with $\pi(y)\ll y/\log y$ and
$\log P_y=\vartheta(y)\ll y$ by Chebyshev's elementary bounds.
Every exceptional prime larger than $y$ satisfies $S_y=1$.

\begin{lemma}[conditioned c-free void; notes Lemma 39.5]\label{lem:void}
For fixed $\kappa,D$ there are $B_0>0$ and $c_v>0$ such that, for every
fixed $B\geq B_0$ and all $X\geq X_v(\kappa,D,B)$,
\begin{equation}\label{eq:void}
 \Prob(H_X=0\mid S_y=1)\leq\e^{-c_vt^3},\qquad y=Bt^3.
\end{equation}
One may choose $c_v$ depending only on $\kappa,D$, uniformly in fixed
$B\geq B_0$ after allowing the threshold to depend on $B$.
\end{lemma}
\begin{proof}
Reveal $c\pmod{L_K}$ and retain the selector condition.  Put
\begin{equation}\label{eq:Z}
 Z(c)=\sum_{\substack{y<p\leq K,\ p\mid L_K\\p\mid c}}\frac1p.
\end{equation}
The restriction $p\mid L_K$ is important: otherwise divisibility by $p$
is not defined on the residue $c$.  In particular not every prime
$p\equiv3\pmod4$ up to $K$ divides an admissible multiplier.
For supported $p>y$ the events $p\mid c$ are independent with probabilities
$1/p$, also after conditioning on $S_y=1$.
Since $0<y/p<1$ and $\e^u-1\leq(\e-1)u$ on $[0,1]$,
\begin{align}
 \E(\e^{yZ}\mid S_y=1)
 &=\prod_{y<p\leq K,\ p\mid L_K}
           \left(1+\frac{\e^{y/p}-1}{p}\right)\notag\\
 &\leq\exp\left\{(\e-1)y\sum_{p>y}p^{-2}\right\}\leq\e^{C_Z},
                                                        \label{eq:Zmoment}\\
 \Prob(Z>\eta\mid S_y=1)&\leq\e^{-\eta y+C_Z}.
                                                        \label{eq:Ztail}
\end{align}
Here $C_Z$ is absolute; even summing over all integers $>y$ bounds the
exponent by an absolute constant.

Choose absolute $a_0>0$, $K_1$ such that
$h(\K(K))\geq a_0\log K$ for $K\geq K_1$ by Lemma~\ref{lem:h}.
The selector rules out every supported $p\leq y$ dividing $c$.
By \eqref{eq:hremove} and a union bound over the remaining prime divisors,
\[
 h(\J_c)\geq h(\K(K))-
       \sum_{y<p\leq K,\ p\mid L_K,\ p\mid c}
          \sum_{k\leq K,\ p\mid k}\frac{\varphi(k)}{k^2}
 \geq (a_0-2Z(c))\log K
\]
for $K$ large.  Fix $0<\eta\leq a_0/4$.  On $Z(c)\leq\eta$ this gives
$h(\J_c)\geq(a_0/2)\log K\geq a_0\kappa t/4$.
Theorem~\ref{thm:harvest} applies to this arbitrary subfamily because
$1\in\J_c$ and $(c,L_{\J_c})=1$; no reducedness modulo all of $L_K$
is being assumed.

Conditional independence and distinctness now give the exact product for
the unpruned active family, followed by its pruned lower estimate:
\begin{align}
 \Prob(H_X=0\mid c,S_y=1)
 &=\prod_\ell(1-f_c(\ell)/\ell)\notag\\
 &\leq\exp\{-\mu_c\}
 \leq\exp\{-a_ht^2h(\J_c)\}
 \leq\e^{-a_ht^3a_0\kappa/4}                 \label{eq:fibrevoid}
\end{align}
on the good event.  The equality remains valid if the selector has prime
coordinates not in $L_K$, as all of them are below $X^{1/2}$ and independent
of the $\ell$-coordinates.  Put $a_v=a_ha_0\kappa/4$ and choose
$B_0$ with $\eta B_0\geq2a_v$.  Averaging \eqref{eq:fibrevoid} and
using \eqref{eq:Ztail} gives
$\Prob(H_X=0\mid S_y=1)\leq\e^{-a_vt^3}+\e^{-2a_vt^3+C_Z}$.
For sufficiently large $X$, this is at most $\e^{-a_vt^3/2}$.
Take $c_v=a_v/2$.
\end{proof}

The role of the selector is quantitative.  A merely Markov bound for the
lost multiplier mass would be far too weak.  The exponential parameter
$y\asymp t^3$ in \eqref{eq:Zmoment} makes the bad fibres as rare as the
required saving, while its own congruence expansion still has acceptable
size.  Revealing $c$ above is not an expansion into the $L_K$ possible
integer progressions; that distinction is crucial in the next section.

\section{Bonferroni assembly on the integers}\label{sec:assembly}
\begin{lemma}[an even Bonferroni majorant]\label{lem:Bonferroni}
For every even integer $r\geq0$ and every integer $h\geq0$, let
$Q_r(h)=\sum_{j=0}^r(-1)^j\binom hj$, with $\binom hj=0$ for $j>h$.
Then
\begin{equation}\label{eq:Bonferroni}
 \begin{gathered}
 Q_r(0)=1,\quad Q_r(h)=\binom{h-1}r\geq0\quad(h\geq1),\\
 \ind_{h=0}\leq Q_r(h)\leq\ind_{h=0}+\binom h{r+1}.
 \end{gathered}
\end{equation}
\end{lemma}
\begin{proof}
Pascal's identity telescopes the alternating sum to
$(-1)^r\binom{h-1}r$ for $h\geq1$.
If $h\leq r$ this is zero; if $h\geq r+1$, then
$\binom h{r+1}=h\binom{h-1}r/(r+1)\geq\binom{h-1}r$.
The case $h=0$ is immediate.
\end{proof}

\begin{theorem}[critical-window assembly; notes Theorem 39.6]\label{thm:assembly}
Fix $\kappa,D$ and $B\geq B_0$ as above.  There are positive constants
$c_a,D_B,C_0$ and $X_a$, depending only on these fixed choices, such that
for every $X\geq X_a$, if $r$ is the least even integer at least $D_Bt^3$,
then
\begin{equation}\label{eq:majorant}
 \nu_X(n)=S_y(n)Q_r(H_X(n))\geq0,\qquad
 \E_{\rm CRT}\nu_X\leq\e^{-c_at^3}.
\end{equation}
It is at least one on every exceptional prime $n>\max(K,y)$.
Its expansion into plain congruence classes has degree at most
$r+\pi(y)$, maximum modulus $q_{\max}$ and total absolute coefficient
sum $T_{\rm abs}$ satisfying
\begin{equation}\label{eq:ledger}
 \begin{aligned}
 r+\pi(y)&=O_{B,D_B}(t^3),\\
 q_{\max}&\leq P_y(KX)^r,\\
 \log q_{\max}&\leq C_\vartheta Bt^3+r(1+\kappa)t,\\
 T_{\rm abs}&\leq2^{\pi(y)}\sum_{j=0}^r\binom{|\A_X|}j,
 \qquad\log T_{\rm abs}\leq C_Lt^4.
 \end{aligned}
\end{equation}
Here $C_\vartheta$ is an absolute Chebyshev constant and
$C_L=C_L(\kappa,B,D_B)$.  For every real $N$ with $\log N\geq C_0t^4$,
\begin{equation}\label{eq:integermean}
 \sum_{1\leq n\leq N}\nu_X(n)\ll_{\kappa,D,B,D_B}N\e^{-c_at^3}.
\end{equation}
\end{theorem}
\begin{proof}
By Lemma~\ref{lem:Bonferroni}, Lemma~\ref{lem:void}, and
Theorem~\ref{thm:moments},
\begin{align}
 \E(Q_r(H_X)\mid S_y=1)
 &\leq\e^{-c_vt^3}+
             \frac{\E((H_X)_{r+1}\mid S_y=1)}{(r+1)!}\notag\\
 &\leq\e^{-c_vt^3}+
             \left(\frac{\e C_Ft^3}{r+1}\right)^{r+1}.
                                                    \label{eq:Btail}
\end{align}
Here $m!\geq(m/\e)^m$ for all integers $m\geq1$.
Choose $D_B>\e C_F$ so large that
$D_B\log(D_B/(\e C_F))\geq2c_v$.  As $r+1\geq D_Bt^3$, the last
term is at most $\e^{-2c_vt^3}$.  Increasing the threshold and setting
$c_a=c_v/2$ gives \eqref{eq:majorant}, after multiplication by
$\Prob(S_y=1)\leq1$.  Positivity, and its value at an exceptional prime
larger than $y$, follow from the identity lemma and \eqref{eq:Bonferroni}.

We now account for \emph{every} term, including the selector.  The identity
\[
 \binom{H_X(n)}j=\sum_{\substack{\mathcal B\subseteq\A_X\\|\mathcal B|=j}}
                              \prod_{A\in\mathcal B}\ind_{E_A}(n)
\]
holds for all integers $n$.  Expanding this and \eqref{eq:selector}
produces terms with coefficients $(-1)^j\mu(d)$ for $d\mid P_y$.
Each intersection is empty or one plain congruence class, by CRT,
with modulus
\begin{equation}\label{eq:termq}
 q=\lcm(d,k_1\ell_1,\ldots,k_j\ell_j)
             \leq P_y(KX)^j.
\end{equation}
In fact any selector divisor sharing a prime with an atom multiplier
makes the term empty, because atom residues are units; ignoring that
cancellation only enlarges our upper ledger.
Every nonempty term modulus divides $\mathcal M$ in \eqref{eq:space}.

For a completely crude inventory, choose $k\leq K$, $\ell\leq X$ and
$u,v\leq X^{1/6}$.  There are at most
\begin{equation}\label{eq:inventory}
 |\A_X|\leq KX^{4/3}\leq\exp\{(4/3+\kappa)t\}
\end{equation}
possibilities.  Conditions of primality, coprimality and divisibility
only reduce this bound.  Therefore
\begin{align}
 \log T_{\rm abs}
 &\leq\pi(y)\log2+\log(r+1)
               +r\log\max(1,|\A_X|)\notag\\
 &\leq C_\pi Bt^3\log2+\log(r+1)
                         +r(4/3+\kappa)t\leq C_Lt^4.\label{eq:termcount}
\end{align}
An absolute $C_\pi$ with $\pi(y)\leq C_\pi y$ suffices in this line;
the sharper $\pi(y)\ll y/\log y$ gives the stated selector degree.
The exponent is $O(rt)$, not $O(r)$; we retain this loss explicitly.
Equation~\eqref{eq:termq} and Chebyshev give the modulus ledger.
Since $r\leq D_Bt^3+2$, all bounds in \eqref{eq:ledger} follow with fixed
constants of the dependencies stated.

For every residue class $a\pmod q$, without any primality restriction,
\[
 \#\{1\leq n\leq N:n\equiv a\pmod q\}=N/q+O(1),
\]
with absolute error at most one.  Summing with the signed coefficients,
and bounding rounding errors by their absolute sum, yields
\begin{equation}\label{eq:transfer}
 \sum_{n\leq N}\nu_X(n)
 =N\E_{\rm CRT}\nu_X+O(T_{\rm abs})
 \leq N\e^{-c_at^3}+\e^{C_Lt^4}.
\end{equation}
Choose $C_0$ greater than twice $C_L$ and twice a fixed upper constant
for $\log q_{\max}/t^4$.  Then $\log N\geq C_0t^4$ ensures both
$T_{\rm abs}\leq N^{1/2}$ and $q_{\max}\leq N^{1/2}$.
The latter is convenient but not necessary for the exact $O(1)$ count.
For large $X$, $N^{1/2}\leq N\e^{-c_at^3}$, since
$\tfrac12\log N\geq(C_0/2)t^4\geq c_at^3$.
This proves \eqref{eq:integermean}.
\end{proof}

This theorem does not establish the literal hypothesis
$H_{\rm PF}'(k\ell;{\rm good})$ of notes \S34.4.  That hypothesis demands
coefficient sum $\e^{O(r)}$ and a majorant of the larger,
$c$-dependently pruned avoider predicate.  Here the family is fixed and
unpruned, its avoider is smaller but still contains every exception, and
the ledger is honestly $\e^{O(rt)}$.  The entire $L_K$-fibre partition
is absent from \eqref{eq:termcount}.  No assertion about a prime-counting
mean for $\nu_X$ is required.

\section{Exceptional primes and all denominators}\label{sec:all}
\begin{proof}[Proof of Theorem~\ref{thm:main} for primes]
For large $N$, put $L=\log N$, take
$t=\alpha L^{1/4}$ and $X=\e^t$, with fixed
$0<\alpha\leq C_0^{-1/4}$.  Every exceptional prime larger than
$\max(K,y)$ is counted with weight at least one by $\nu_X$, so
\[
 E_{\rm pr}(N)\leq K+y+\sum_{n\leq N}\nu_X(n)
 \ll K+y+N\exp\{-c_a\alpha^3L^{3/4}\}.
\]
The omitted range is negligible: $K\leq\e^{\kappa\alpha L^{1/4}}$ and
$y=B\alpha^3L^{3/4}$, whereas the logarithm of the last main bound is
$L-c_a\alpha^3L^{3/4}$.  Reducing the saving constant if needed gives
$E_{\rm pr}(N)\ll N\e^{-c_p(\log N)^{3/4}}$ for a fixed $c_p>0$.
Increasing the outer constant covers bounded $N\geq3$.
\end{proof}

\begin{proof}[Transfer to all denominators; {\cite[Section 6]{LL}}]
If $p\mid n$ and $4/p=\sum_i1/x_i$, then
$4/n=\sum_i1/((n/p)x_i)$.  Every prime factor of an exceptional integer
is therefore exceptional.  Thus exceptional integers lie in the
multiplicative semigroup generated by exceptional primes, including the
empty product 1.  This containment, not its converse, is all that is used.

Put $g(u)=u^{3/4}$.  Choose $c_0>0$ and $x_0\geq\e$ such that for every
$x\geq x_0$ the prime bound reads
$E_{\rm pr}(x)\leq x\e^{-c_0g(\log x)}$; a decrease of the saving
constant absorbs its prefactor.  Fix $0<\eta_0<c_0$, let $L=\log x$,
and put $\delta=\eta_0g(L)/L=\eta_0L^{-1/4}$, which is less than
$1/4$ for large $x$.  Rankin's inequality gives
\begin{equation}\label{eq:Rankin}
 E(x)\leq x^{1-\delta}
       \prod_{p\leq x,\ p\ \mathrm{exceptional}}(1-p^{-1+\delta})^{-1}.
\end{equation}
The higher prime-power terms in the logarithm of the product are
$O(\sum_p p^{-3/2})=O(1)$ uniformly for $\delta<1/4$.
The finitely many primes below $x_0$ also cost a fixed constant.
Partial summation gives
\[
 \sum_{x_0<p\leq x,\ p\ \mathrm{exceptional}}p^{-1+\delta}
 \ll1+\e^{\delta L-c_0g(L)}+
             \int_{\log x_0}^{L}\e^{\delta u-c_0g(u)}\,du.
\]
Since $g(u)/u=u^{-1/4}$ decreases, $\delta u\leq\eta_0g(u)$ for
$u\leq L$.  The endpoint is bounded by one and the integral by
$\int_{\log x_0}^{\infty}\e^{-(c_0-\eta_0)u^{3/4}}\,du<\infty$.
The Euler product in \eqref{eq:Rankin} is bounded uniformly in $x$.
Consequently $E(x)\ll x^{1-\delta}=x\e^{-\eta_0(\log x)^{3/4}}$.
After fixing $\kappa$ absolutely, all constants are absolute, completing
the claimed argument.
\end{proof}

\section{The heuristic mass-driven three-quarter ceiling}\label{sec:ceiling}
We record the large-sieve convention of \cite[Section~5]{LL}, even though
our interval transfer is Bonferroni rather than an invocation of that sieve.
Montgomery \cite[\S8, p.~561]{Montgomery} states, for an integer set in an
interval of length $Y$ avoiding $0\leq\omega(\ell)<\ell$ classes at each
prime in a finite set $\mathcal P$,
\[
 |\mathcal N|\leq\frac{Y+Q^2}{S(Q)},\qquad
 S(Q)=\sum_{s\leq Q,\ s\mid\prod\mathcal P}\mu^2(s)
                   \prod_{\ell\mid s}\frac{\omega(\ell)}{\ell-\omega(\ell)}.
\]
This is the large sieve, not the distinct larger sieve.  Writing
$g(\ell)=\omega(\ell)/(\ell-\omega(\ell))$ and
$G=\prod(1+g(\ell))$, the Rankin tail with $v=1/\log X$ gives,
when $\ell\leq X$ and $\omega(\ell)\leq\ell/2$,
\[
 \frac{G-S(Q)}G\leq
 \exp\left\{-\frac{\log Q}{\log X}
                     +2(\e-1)\sum_\ell\frac{\omega(\ell)}\ell\right\}.
\]
Thus mass $\mu\asymp t^{\mathcal B}$ and this product-tail budget balance
$t\mu\lesssim\log N$.  Our Bonferroni ledger has the same balance:
$r\asymp\mu$ and $\log T_{\rm abs}=O(rt)$.
The resulting exponent is $\theta=\mathcal B/(\mathcal B+1)$.
Vaughan's prime-only supply has $\mathcal B=2$; here $\mathcal B=3$.
(The symbol $\mathcal B$ in this paragraph is not the selector constant $B$.)

For completeness, the cubic total identity-class supply asserted in notes
\S18 has a short elementary verification.  For $M\equiv3\pmod4$, put
$A=(M+1)/4$ and let $F(M)$ count the distinct multiplier-identity classes
modulo $M$, using every factorization $A=uvw$ once as a source and then
deduplicating residues.  They are exactly $-4D\pmod M$ with $D\mid A^2$:
one direction uses $D=u^2w$; conversely at $p^e\Vert A$, $p^d\Vert D$,
set the exponents in $u,w,v$ to
$\lfloor d/2\rfloor$, $d-2\lfloor d/2\rfloor$,
$e-\lfloor d/2\rfloor-(d-2\lfloor d/2\rfloor)$.
The divisors $D<A$ give distinct residues since $0<4D<M$, so
\[
 \tfrac12(\tau(A^2)-1)\leq F(M)\leq\tau(A^2).
\]
The coefficient identity
$\tau(n^2)=\sum_{d^2\mid n}\mu(d)\tau_3(n/d^2)$ follows from
$\sum_{e\geq0}(2e+1)z^e=(1-z^2)/(1-z)^3$.
Also $\sum_{m\leq x}\tau_3(m)=\frac12x\log^2x+O(x\log x)$:
sum $\lfloor x/(ab)\rfloor$ over $ab\leq x$, whose total rounding
error is $O(x\log x)$, and use the elementary harmonic sums for
$\sum_{ab\leq x}1/(ab)=\frac12\log^2x+O(\log x)$.
Summing over $d$ and then partially summing yields
\[
 \sum_{n\leq x}\frac{\tau(n^2)}n
     =\frac{\log^3x}{6\zeta(2)}+O(\log^2x),\qquad
 \sum_{M\leq Q,\ M\equiv3\ (4)}\frac{F(M)}M\asymp\log^3Q.
\]
Changing $1/(4A-1)$ to $1/(4A)$ costs only $O(1)$.
Thus the entire identity-class supply, with distinct moduli and residues,
already has exponent three; counting the same $M$ again through another
factorization $M=k\ell$ does not create new supply.

\begin{remark}[scope of the ceiling]
The mass-and-budget calculation suggests that changing only a fixed level
of distribution does not change the predicted exponent $3/4$ when the
saving comes only from this summed identity-class mass and the truncation
retains the stated $t\mu$ budget.  This is a \emph{heuristic ceiling}, not
an optimality theorem: a sufficient product-tail or coefficient budget has
not been proved necessary for every implementation of the method.
In particular it does not rule out other congruence-and-sieve arguments,
extra identities, or use of arithmetic arrangement beyond total mass.  The
stronger universal formulation does not follow from the supply calculation;
notes \S18
explicitly makes this same restriction.  Nothing about the ceiling is
used in the proof of Theorem~\ref{thm:main}.
\end{remark}

\section{The three external-referee checkpoints}\label{sec:checkpoints}
The following are exactly the three mathematical checkpoints identified
in notes \S39.7, with their locations in this standalone presentation.
Internal re-derivation is not external refereeing.  Each checkpoint has
now been re-derived independently, from the statements rather than by
transcription (\cite[\S76.4]{Notes}); the findings are recorded after
each item.  The theorem carries the label \textbf{internally proved;
internal checks only, not externally refereed}, with the internal
record described in Theorem~\ref{thm:main}.

\begin{enumerate}
\item \emph{Residue-resolved local factors.}
Re-derive \eqref{eq:residuebox} (notes (39.11)) uniformly for every
$k\leq X^\kappa$, every divisor $g\mid k$, and every unit residue.
The $q_0=\lcm(g,\rad k)$ class count, Shiu interval hypotheses, and the
factor $\exp\{-2\sum_{p\mid k}(p-1)^{-1}\}$ must all be retained.
The upper bound uses BT, not a BV main term below its error.
Lemma~\ref{lem:profile} gives the full calculation and separately checks
that its fixed-$k$ lower main term does dominate the BV error.
\emph{Re-derivation:} confirmed.  The checkpoint is not load-bearing:
it enters only the ordered-atom replay of Theorem~\ref{thm:moments},
whose conditional-independence proof needs only
Corollary~\ref{cor:fibremass}.

\item \emph{Pruned supply for the data-dependent subfamily.}
Check Theorem~\ref{thm:harvest} (notes Theorem~34.8) and its use with
$\J_c$ in Corollary~\ref{cor:fibremass} and Lemma~\ref{lem:void}
(notes Lemma~39.5).  This includes the power-range lattice and Shiu
estimates, every second-moment boundary error, $1\in\J_c$ for uniform
pruning, the polylogarithmic modulus multiplicity after pruning, and the
fact that the lower mass occurs in the \emph{canonical} $z=x^{1/6}$ boxes.
These boxes, not a larger per-prime cutoff, define our fixed atom family.
The exponential bad-fibre calculation uses only primes supported by $L_K$.
\emph{Re-derivation:} confirmed line by line.  The pruned-supply
component is replaced by Theorem~\ref{thm:unpruned}, i.e.\ by
Lemma~\ref{lem:wsecond} (one Shiu application in progressions modulo
$d\leq K^2$ in boxes above $K^{10}$, plus the elementary modulus-one
case), Cauchy--Schwarz, and Bombieri--Vinogradov with a fixed $R$; the
elementary lattice estimates, the $\J_c$ quantifiers and fibre bounds,
and the selector's exponentially small bad-fibre estimate remain
required.

\item \emph{Prime-power consistency versus lcm collapse.}
Replay \eqref{eq:localrelative}--\eqref{eq:cancel} and the ordered moment
induction (notes (39.14)--(39.19)), including unequal prime-power exponents.
A shared conditioned $p^e$ contributes $\varphi(p^e)$, an unconditioned one
$p^e$; the residue share cancels $\varphi(p^e)$, leaving only $p/(p-1)$
in the latter case.  Any missed factor would be raised through
$r\asymp t^3$.  The independent-Bernoulli proof of
Theorem~\ref{thm:moments} provides a further check, using the uniform
all-fibre upper bound rather than a Poisson approximation.
\emph{Re-derivation:} confirmed, including the exact conditional
independence on $\mathbb Z/\mathcal M$, the distinctness of atoms at a
common $\ell$, the exact falling-factorial identity, the local
prime-power table, and the $\e^{O(t^4)}$ ledger.
\end{enumerate}

Two further items outside the list were checked in the re-derivation:
the Rankin transfer of Section~\ref{sec:all} and the parameter
inequalities \eqref{eq:floor} from $\kappa<1/240$.  Both are correct.

No step in the displayed implication is supplied by the unproved literal
$H_{k\mathrm{BV}}$ or by a partition-free sieve hypothesis, and none invokes
Elliott--Halberstam.  The proof calculations above have been checked in
preparing this condensation; no unresolved step has been silently replaced
by an assumption.  Nevertheless an external expert's reading of these checks and a
primary-literature priority audit are still required before treating
the theorem as externally established.  The stronger universal sieve
ceiling and an exact finite-truncation interpretation of the proxy Euler
share are not asserted.

\begin{thebibliography}{99}
\bibitem[D]{Davenport}
H. Davenport, \emph{Multiplicative Number Theory}, third edition,
revised by H. L. Montgomery, Graduate Texts in Mathematics 74,
Springer, 2000, Chapter 28.
\bibitem[LL]{LL}
\emph{A logarithmic improvement of Vaughan's bound for the
Erd\H{o}s--Straus exceptional set}, internally reviewed standalone note,
\texttt{paper/vaughan-loglog-note.tex}, this repository, 2026.
\bibitem[M]{Montgomery}
H. L. Montgomery, The analytic principle of the large sieve,
\emph{Bull. Amer. Math. Soc.} \textbf{84} (1978), 547--567.
\bibitem[MV]{MV}
H. L. Montgomery and R. C. Vaughan, The large sieve,
\emph{Mathematika} \textbf{20} (1973), 119--134.
\bibitem[N]{Notes}
\emph{Erd\H{o}s--Straus research notebook}, \texttt{notes.md}, this
repository, \S\S16,18,34,39,41,76; internally reviewed working notes,
2026.
\bibitem[S]{Shiu}
P. Shiu, A Brun--Titchmarsh theorem for multiplicative functions,
\emph{J. Reine Angew. Math.} \textbf{313} (1980), 161--170.
\bibitem[V]{Vaughan}
R. C. Vaughan, On a problem of Erd\H{o}s, Straus and Schinzel,
\emph{Mathematika} \textbf{17} (1970), 193--198.
\end{thebibliography}
\end{document}
